\begin{lemma}[Starred-prefix preparation for Algorithm 5.6]
Fix a finite loopless multigraph $G$, palette $[k]$, and root edge $e_0$.
Let $\mathcal E$ be a positive-rung ETT with input coloring $\varphi_{\rm in}$,
displayed tree $T$, generating coloring $\varphi$, rung $n$, maximum property,
and base $B$. Assume a \noderef{TechnicalInductionContext} at bound $n$,
$\Delta(G)+1\le k$, and $11\le t(G)$. Assume the following uniform bound:
for every arbitrary-input maximum ETT $\mathcal L$ with rung $n$, generating
coloring $\psi$, base $C$, and good hierarchy $K$, its terminal ladder tree
$L_n$ satisfies $2n+11\le |\overline\psi(L_n)|$.

Let $U$ have a corresponding certificate $C_U$ relative to $\mathcal E$
and $\varphi$, maximum property, and a base $B_U$ with
$B_U.\mathit{core}=B.\mathit{core}$. Let $H$ be a hierarchy on this actual
replacement ETT, with $d$ dividers and good witness $\mathcal G$.
Assume that $|B.\mathit{core}|<|U|$ when the original terminal kind is PE,
and $|T_n|<|U|$ when it is not PE. Here tree lengths count displayed edges.
Set $A=\noderef{HierarchyStarredStage}(H,d)$ and
$D^+=\noderef{HierarchyTerminalStoredColors}(H)$.
Then $\varphi$ is proper and leaves exactly $e_0$ uncolored,
$V(A)$ and $V(U)$ are elementary under $\varphi$,
$A$ is an exact ordered initial segment of $U$ with $|A|<|U|$,
$|V(A)|$ is odd, and
\[
 \left|\bigcup_{\theta\in D^+}\mathcal G.\Gamma^d_\theta\right|
       < |\overline\varphi(U)|.
\]
No terminal minus-closedness or nonclosedness hypothesis on $U$ is needed.
\end{lemma}
\begin{proof}
Write $\mathcal L=C_U.\mathit{replacementETT}$.
The rung agreement and series-presentation fields of
\noderef{CorrespondingExtendedTashkinovTree} identify its rung with $n$.
They also identify every positive ladder tree, terminal kind, and raw
stored set with the corresponding original data. All uses of these
identifications below are through those fields, with finite indices
transported by their numerical values.

First inspect the terminal certificate of the original
\noderef{TashkinovSeries}. Since $n\ge1$ by
\noderef{ExtendedTashkinovTree}, it cannot be the initial constructor.
In the SE constructor of \noderef{TashkinovSeriesStep} the selected
coloring is proper and leaves exactly $e_0$ uncolored by its
\noderef{MaximumDefectiveChoice} certificate; its current-color
equation identifies this coloring with the terminal coloring.
In PE these two laws are explicit current-color laws. In RE the current
coloring equals the preceding coloring, which has these two laws.
The generating-color equation of the ETT identifies the terminal
coloring with $\varphi$. This proves the first two assertions. The same
argument applies to any positive-rung certified ETT used below.

For any maximum ETT of rung $n$, assemble the six strict lower-rung
assertions from the six lower fields of the assumed
\noderef{TechnicalInductionContext}, in their listed order. Assemble
the four current assertions (iii)--(vi) from its four current fields.
For $\mathcal L$, its rung equality with $n$ changes every strict bound
and every current-rung equality into the corresponding context bound.
Apply \noderef{GoodHierarchyElementaryLarge} to $\mathcal L$, its
given maximum property, $B_U,H,\mathcal G$, these two assembled
packages, $11\le t(G)$, and the uniform missing bound. The latter is
instantiated on this actual replacement ETT and this actual good
hierarchy. It follows that $V(U)$ is elementary.

We will also use pre-terminal elementarity on the ladder tree $L_n$ of
$\mathcal L$. If $n=1$, the elementarity conclusion of
\noderef{TechnicalTheoremBase}, applied to its initial certified
series, makes $L_1$ elementary under the slot-zero coloring.
If $n>1$, \noderef{CertifiedETTPredecessor} supplies an actual maximum
ETT of rung $n-1$, whose displayed tree is exactly $L_n$ and whose
generating coloring is the slot-$(n-1)$ coloring. Apply the context's
lowerOne field to this actual certificate; $0<n-1<n$ verifies both
bounds. Thus in both cases $L_n$ is elementary under the pre-step
coloring. This argument never invokes a zero-rung ETT.

Suppose first $d=0$. By \noderef{HierarchyStarredStage}, $A=B_U.\mathit{core}$.
If the terminal kind is PE, apply
\noderef{ParallelClosureUnionElementary} to $\mathcal L$, its maximum
property, the assembled six lower assertions, and the coloring $\varphi$.
Its properness and unique-uncoloredness were checked above; its stability
relative to itself follows from reflexivity in \noderef{sc2}.
The closure and ordered-union data are precisely the parallel field of
\noderef{HierarchyBase}. The conclusion makes $V(A)$ elementary.

In this PE case the parallelBase field of \noderef{GoodHierarchy}
also supplies at least four colors in
\[
 (\overline\varphi(L_n)\cap\overline\varphi(B_U.\mathit{closure}))
  -\bigcup_{\theta\in D_{n,0}}\mathcal G.\Gamma^0_\theta.
\]
Choose any one, $a$. Apply \noderef{ParallelCommonMissingClosed} to
the terminal PE, this base, and the pre-step elementarity just proved.
It says that every common missing color in the displayed intersection
is closed on the core. Moreover the ordered-union clause of
\noderef{OrderedTreeUnion} puts $L_n$ inside the core, so $a$ is missing
on $A$ as well. Thus $a$ is both missing on and closed at $A$.

If the terminal kind is not PE, the positive terminal certificate is
SE or RE. By the nonparallel field of \noderef{HierarchyBase},
$A=L_n$. In SE, the selected coloring is stable on $L_n$ relative to
the pre-step coloring; in RE the two colorings are equal. The vertex
clause of \noderef{StableColoring} therefore preserves every missing
set on $L_n$ and transports the pre-step elementarity to $\varphi$.
Both constructors say that $V(L_n)$ was closed under the pre-step
coloring in the sense of \noderef{ColoredVertexSetClosed}. It remains
closed under $\varphi$: a hypothetical $\varphi$-colored boundary
edge in a missing color has a tree endpoint by \noderef{BoundaryEdges}
and is therefore incident with the tree by \noderef{TreeIncidentEdge};
symmetry of the stable
relation from \noderef{sc2}, followed by its incident-edge clause,
gives that same color on that exact edge under the pre-step coloring.
The vertex clause supplies missingness there under the pre-step
coloring, contradicting its closedness. In RE this is simply equality
of colorings. Pick a vertex of $A$ (its vertex array has positive
length) and a missing color there by
\noderef{MissingColorsNonemptyOfMaximumDegree}. The closedness just
proved makes this color closed on $A$ under \noderef{TreeColorSetClosed}.

Now suppose $d>0$. Set $r=\operatorname{last}(d)$, so $r\ge1$.
The definition of \noderef{HierarchyStarredStage} gives
$A=H.\mathit{stage}(r)$.
Apply \noderef{GoodHierarchyPrefixRepackage} to $\mathcal L$, its
maximum property, $B_U,H,\mathcal G$, and this positive index. It
returns an actual maximum ETT of rung $n$ whose displayed tree is
exactly $A$, with a base, a truncated good hierarchy, and the original
series up to heterogeneous equality. Applying
\noderef{GoodHierarchyElementaryLarge} to this actual package makes
$V(A)$ elementary: rung equality transports the lower and current
context packages, and the uniform missing hypothesis applies to
this prefix ETT and its returned good hierarchy. The same prefix
theorem also gives minus-closedness of $A$ with respect to
\[
 R_A=\bigcup_{\theta\in D_{n,d}}\mathcal G.\Gamma^{d-1}_\theta.
\]
By \noderef{HierarchyStoredColorsAntitone},
$D_{n,d}\subseteq D_{n,d-1}$, so each pair in this union has exactly
two colors by the twoColors field of \noderef{GoodHierarchy}.
By \noderef{HierarchyStoredColorsCardinality}, $|D_{n,d}|\le n$.
Counting a union by the sum of its member cardinalities gives
$|R_A|\le2n$.
The ladder tree $L_n=H.\mathit{stage}(0)$ is an exact prefix of $A$:
compose the nested fields of \noderef{TashkinovHierarchy} through
stage $d$. Hence the uniform bound on $\mathcal L$ gives
$2n+11\le|\overline\varphi(A)|$ by the union definition
\noderef{TreeMissingColors}. Choose a missing color outside $R_A$;
the returned minus-closedness, by \noderef{TreeColorSetMinusClosed},
makes it closed on $A$.

In all cases $A$ is elementary and has a closed missing color.
Apply \noderef{ElementaryClosedMissingColorOdd} with the properness
proved above to conclude that $|V(A)|$ is odd.

For exact prefix and strictness, when $d>0$ the last nested and strict
fields of \noderef{TashkinovHierarchy}, together with terminalStage,
give $A\preccurlyeq U$ and $|A|<|U|$. When $d=0$, coreInsideFirst and
terminalStage give $A\preccurlyeq U$. In PE the assumed strict bound
and $B_U.\mathit{core}=B.\mathit{core}$ give strictness. In SE/RE,
$A=L_n=T_n$ by nonparallel and the positive ladder agreement, and
the other assumed strict bound gives strictness.

Finally \noderef{HierarchyTerminalStoredColorsSubsetLast} gives
$D^+\subseteq D_{n,d}$. Every old last pair belonging to an owner
in $D^+$ consequently has cardinality two. The indexed stored-set
bound \noderef{HierarchyStoredColorsCardinality} and rung agreement
give $|D^+|\le n$. Therefore its reserved union has cardinality at
most $2n$. The abovePenultimate field of
\noderef{ExtendedTashkinovTree} puts $L_n$ inside $U$; explicitly,
\noderef{TreeSequenceInitialSegmentVertexSetSubset} carries each
missing witness on $L_n$ into $U$. The uniform bound applied to
$\mathcal L,B_U,H,\mathcal G$ gives
$2n+11\le|\overline\varphi(U)|$. Comparing these two inequalities
proves the required strict cardinal gap.
\end{proof}
